\chapter{Quick Start}

\section{What is OpenClaw}

\begin{quote}
\textit{"De-claw! De-claw!"} — Probably said by a space lobster
\end{quote}

OpenClaw is an AI Agent Gateway that supports multiple chat platforms including WhatsApp, Telegram, Discord, iMessage, and more.
It connects chat applications to AI coding agents through a single Gateway process, supporting both local and remote deployment.

\subsection{How It Works}

\begin{center}
\begin{tikzpicture}[
    node distance=0.8cm and 1.5cm,
    every node/.style={font=\small},
    box/.style={draw=blue!50, fill=blue!5, rounded corners=4pt, minimum width=2.8cm, minimum height=0.7cm, align=center},
    gw/.style={draw=orange!60, fill=orange!10, rounded corners=6pt, minimum width=3.5cm, minimum height=1cm, align=center, font=\bfseries},
    arr/.style={-{Stealth[length=2.5mm]}, thick, draw=gray!60}
]
    % Left: Chat apps
    \node[box] (wa) {WhatsApp};
    \node[box, below=of wa] (tg) {Telegram};
    \node[box, below=of tg] (dc) {Discord};
    \node[box, below=of dc] (im) {iMessage};

    % Center: Gateway
    \node[gw, right=2cm of $(tg)!0.5!(dc)$] (gw) {Gateway};

    % Right: Outputs
    \node[box, right=2cm of gw, yshift=0.8cm] (agent) {AI Agent};
    \node[box, right=2cm of gw, yshift=-0.8cm] (web) {Web UI};

    % Arrows
    \foreach \src in {wa, tg, dc, im} {
        \draw[arr] (\src) -- (gw);
    }
    \draw[arr] (gw) -- (agent);
    \draw[arr] (gw) -- (web);
\end{tikzpicture}
\end{center}

The Gateway is the \textbf{single source of truth} for sessions, routing, and channel connections.

\section{Three-Step Installation}

\subsection{Step 1: Install OpenClaw}

\begin{lstlisting}[language=bash]
npm install -g openclaw@latest
\end{lstlisting}

\subsection{Step 2: Run the Onboarding Wizard}

\begin{lstlisting}[language=bash]
openclaw onboard --install-daemon
\end{lstlisting}

The onboarding wizard automatically configures:
\begin{itemize}
  \item Model/authentication (OAuth or API key recommended)
  \item Channel login (WhatsApp, Telegram, etc.)
  \item Daemon (systemd/launchd auto-start on boot)
\end{itemize}

\subsection{Step 3: Start the Gateway}

\begin{lstlisting}[language=bash]
openclaw channels login
openclaw gateway --port 18789
\end{lstlisting}

Open your browser and visit \texttt{http://127.0.0.1:18789/} to start chatting.

\section{Quickest Experience: Control UI}

No channel configuration needed --- chat directly in the browser:

\begin{lstlisting}[language=bash]
openclaw dashboard
\end{lstlisting}

Open \texttt{http://127.0.0.1:18789/} in your browser to start talking with the agent.

\section{Authentication Configuration}

\subsection{Anthropic API Key (Recommended)}

\begin{lstlisting}[language=bash]
export ANTHROPIC_API_KEY="..."
openclaw models status
\end{lstlisting}

Or store the key in the daemon configuration:

\begin{lstlisting}[language=bash]
cat >> ~/.openclaw/.env <<'EOF'
ANTHROPIC_API_KEY=...
EOF
\end{lstlisting}

\subsection{Verify Installation}

\begin{lstlisting}[language=bash]
openclaw models status
openclaw doctor
\end{lstlisting}

\section{Core Concepts at a Glance}

OpenClaw consists of the following core components, each explained in detail in subsequent chapters:

\begin{center}
\begin{tabular}{lll}
\toprule
\textbf{Concept} & \textbf{Description} & \textbf{Config Path} \\
\midrule
Gateway & Core AI gateway, WS/HTTP server & \texttt{gateway.*} \\
Agent & AI assistant instance, bound to models/skills & \texttt{agents.list[]} \\
Channel & Chat platform adapter & \texttt{channels.<name>.*} \\
Binding & User/group to Agent routing & \texttt{agents.list[].bindings[]} \\
Session & A conversation with context history & \texttt{session.*} \\
Plugin & Extension module for tools/channels/hooks & \texttt{plugins.entries.*} \\
Skill & Agent skill (tools + system prompt) & \texttt{skills.entries.*} \\
Memory & Long-term memory, vector + full-text search & \texttt{memorySearch.*} \\
Sandbox & Docker sandbox for isolated execution & \texttt{agents.defaults.sandbox.*} \\
\bottomrule
\end{tabular}
\end{center}

\section{Configuration File}

The configuration file is located at \texttt{\textasciitilde{}/.openclaw/openclaw.json} (JSON5 format).

\begin{lstlisting}[language=json]
{
  channels: {
    whatsapp: {
      allowFrom: ["+15555550123"],
      groups: { "*": { requireMention: true } }
    }
  },
  messages: {
    groupChat: { mentionPatterns: ["@openclaw"] }
  }
}
\end{lstlisting}

If you \textbf{make no changes}, OpenClaw will run using the built-in Pi binary in RPC mode and create independent sessions per sender.
